\begin{producciones}
\regla{Inicial}{\nt{ListaClases} \tl{eof}}
\regla{ListaClases}{\nt{Clase} \nt{ListaClases} \alt{} \nt{Interfaz} \nt{ListaClases} \alt{} \eps}
\regla{Clase}{\tl{class} \tl{idClase} \nt{GenericidadOpcional} \nt{HerenciaOpcional} \tl{\{} \nt{ListaMiembros} \tl{\}}}
\regla{Interfaz}{\tl{interface} \tl{idClase} \nt{GenericidadOpcional} \nt{ExtensionOpcional} \tl{\{} \nt{ListaMetodosInterfaz} \tl{\}}}
\end{producciones}

\noindent{\small\color{grisMedio}Una Clase o Interfaz sigue limitada a UN único parámetro de tipo, también con el logro de genericidad avanzada.\par}

\begin{producciones}
\regla{GenericidadOpcional}{\tl{<} \tl{idGen} \tl{>} \alt{} \eps}
\regla{HerenciaOpcional}{\tl{extends} \nt{TipoReferencia} \alt{} \tl{implements} \nt{TipoReferencia} \alt{} \eps}
\regla{ExtensionOpcional}{\tl{extends} \nt{TipoReferencia} \alt{} \eps}
\regla{ListaMiembros}{\nt{Miembro} \nt{ListaMiembros} \alt{} \eps}
\regla{ListaMetodosInterfaz}{\nt{MetodoInterfaz} \nt{ListaMetodosInterfaz} \alt{} \eps}
\end{producciones}

\noindent{\small\color{grisMedio}La visibilidad se factoriza al frente de los tres tipos de miembro. Eso hace que el \texttt{public} del \nt{Constructor} original quede absorbido por \nt{VisibilidadOpcional}, y que \nt{Constructor} pase a competir con \nt{Tipo} por el token idClase: \texttt{C()\{\}} (constructor) vs \texttt{C x;} (atributo de tipo C).\par}

\begin{producciones}
\regla{Miembro}{\nt{VisibilidadOpcional} \nt{MiembroSinVisibilidad}}\hfill\mbox{\footnotesize\color{grisMedio}[+\,\texttt{visibilidad}]}
\regla{Miembro}{\nt{MiembroSinVisibilidad}}\hfill\mbox{\footnotesize\color{grisMedio}[-\,\texttt{visibilidad}]}
\regla{Miembro}{\nt{Constructor}}\hfill\mbox{\footnotesize\color{grisMedio}[-\,\texttt{visibilidad}]}
\regla{VisibilidadOpcional}{\tl{public} \alt{} \tl{private} \alt{} \eps}\hfill\mbox{\footnotesize\color{grisMedio}[+\,\texttt{visibilidad}]}
\end{producciones}

\noindent{\small\color{grisMedio}Sin el logro, el \texttt{public} sólo aparece en el constructor y es lo que lo distingue del resto de los miembros: es la regla original de la cátedra.\par}

\begin{producciones}
\regla{Constructor}{\tl{public} \tl{idClase} \nt{CuerpoDeMetodo}}\hfill\mbox{\footnotesize\color{grisMedio}[-\,\texttt{visibilidad}]}
\regla{MiembroSinVisibilidad}{\tl{static} \nt{TipoMetodo} \tl{idMetVar} \nt{CuerpoDeMetodo}}
\regla{MiembroSinVisibilidad}{\tl{void} \tl{idMetVar} \nt{CuerpoDeMetodo}}
\regla{MiembroSinVisibilidad}{\nt{TipoPrimitivo} \nt{DimensionesOpcionales} \tl{idMetVar} \nt{AtributoOMetodo}}
\regla{MiembroSinVisibilidad}{\tl{idGen} \nt{DimensionesOpcionales} \tl{idMetVar} \nt{AtributoOMetodo}}
\regla{MiembroSinVisibilidad}{\tl{idClase} \nt{MiembroDesdeIdClase}}
\end{producciones}

\noindent{\small\color{grisMedio}El NT insertado que decide entre constructor y miembro con tipo de referencia. Con visibilidad el constructor ya no lleva \texttt{public} propio, así que compite con el atributo/método de tipo referencia por el token idClase.\par}

\begin{producciones}
\regla{MiembroDesdeIdClase}{\nt{CuerpoDeMetodo}}\hfill\mbox{\footnotesize\color{grisMedio}[+\,\texttt{visibilidad}]}
\regla{MiembroDesdeIdClase}{\nt{TipoGenericoOpcional} \nt{DimensionesOpcionales} \tl{idMetVar} \nt{AtributoOMetodo}}
\end{producciones}

\noindent{\small\color{grisMedio}El NT insertado que decide entre atributo y método.\par}

\begin{producciones}
\regla{AtributoOMetodo}{\nt{InicializacionDeAtributo} \tl{;}}
\regla{AtributoOMetodo}{\nt{CuerpoDeMetodo}}
\end{producciones}

\noindent{\small\color{grisMedio}NT dedicado: la inicialización de un ATRIBUTO es un logro distinto que la de una variable local, y tienen que poder prenderse por separado.\par}

\begin{producciones}
\regla{InicializacionDeAtributo}{\tl{=} \nt{ExpresionCompuesta}}\hfill\mbox{\footnotesize\color{grisMedio}[+\,\texttt{atrinit}]}
\regla{InicializacionDeAtributo}{\eps}
\end{producciones}

\noindent{\small\color{grisMedio}Se comparte a propósito: es la misma construcción en método, método static y constructor.\par}

\begin{producciones}
\regla{CuerpoDeMetodo}{\nt{ArgsFormales} \nt{Bloque}}
\end{producciones}

\noindent{\small\color{grisMedio}Los métodos de interfaz no llevan visibilidad: el logro habla de Métodos, Atributos y Constructores, que son los \nt{Miembro} de una clase.\par}

\begin{producciones}
\regla{MetodoInterfaz}{\nt{TipoMetodo} \tl{idMetVar} \nt{ArgsFormales} \tl{;}}
\regla{TipoMetodo}{\nt{Tipo} \alt{} \tl{void}}
\regla{Tipo}{\nt{TipoBase} \nt{DimensionesOpcionales}}
\regla{TipoBase}{\nt{TipoPrimitivo} \alt{} \nt{TipoReferencia} \alt{} \tl{idGen}}
\regla{DimensionesOpcionales}{\tl{[} \tl{]} \nt{DimensionesOpcionales} \alt{} \eps}
\regla{TipoReferencia}{\tl{idClase} \nt{TipoGenericoOpcional}}
\end{producciones}

\noindent{\small\color{grisMedio}[L:float] — \texttt{float} es token del léxico desde la Etapa 1 pero no estaba en la gramática; sin esta alternativa el logro de Etapa 1 queda inutilizable.\par}

\begin{producciones}
\regla{TipoPrimitivo}{\tl{boolean} \alt{} \tl{char} \alt{} \tl{int}}
\regla{TipoPrimitivo}{\tl{float}}\hfill\mbox{\footnotesize\color{grisMedio}[+\,\texttt{float}]}
\end{producciones}

\noindent{\small\color{grisMedio}En POSICIÓN DE TIPO no se admite el diamante: \texttt{List<> x;} es inválido.\par}

\begin{producciones}
\regla{TipoGenericoOpcional}{\tl{<} \nt{ArgumentoDeTipo} \tl{>} \alt{} \eps}
\end{producciones}

\noindent{\small\color{grisMedio}[L:generico] — la recursión por TipoGenericoOpcional es la que habilita el anidamiento \texttt{List<List<Persona>>}. Sin el logro: \texttt{ArgumentoDeTipo -> idGen | idClase}.\par}

\begin{producciones}
\regla{ArgumentoDeTipo}{\tl{idGen}}
\regla{ArgumentoDeTipo}{\tl{idClase} \nt{TipoGenericoOpcional}}\hfill\mbox{\footnotesize\color{grisMedio}[+\,\texttt{generico}]}
\regla{ArgumentoDeTipo}{\tl{idClase}}\hfill\mbox{\footnotesize\color{grisMedio}[-\,\texttt{generico}]}
\regla{ArgsFormales}{\tl{(} \nt{ListaArgsFormalesOpcional} \tl{)}}
\regla{ListaArgsFormalesOpcional}{\nt{ListaArgsFormales} \alt{} \eps}
\regla{ListaArgsFormales}{\nt{ArgFormal} \nt{RestoListaArgsFormales}}
\regla{RestoListaArgsFormales}{\tl{,} \nt{ArgFormal} \nt{RestoListaArgsFormales} \alt{} \eps}
\regla{ArgFormal}{\nt{Tipo} \tl{idMetVar}}
\end{producciones}

\noindent{\small\color{grisMedio}Con [L:varclasica] una declaración local puede empezar con idClase (\texttt{C x;}), y una sentencia-expresión también (\texttt{C.m();}). Se factoriza ese idClase y el desempate lo hace el token siguiente: \texttt{.} es expresión, cualquier otro es declaración. La rama de expresión reusa \nt{ColaDeExpresionDesdePrimario} para no duplicar la jerarquía de expresiones.\par}

\begin{producciones}
\regla{Bloque}{\tl{\{} \nt{ListaSentencias} \tl{\}}}
\regla{ListaSentencias}{\nt{Sentencia} \nt{ListaSentencias} \alt{} \eps}
\regla{Sentencia}{\tl{;}}
\regla{Sentencia}{\nt{VarLocal} \tl{;}}
\regla{Sentencia}{\nt{Return} \tl{;}}
\regla{Sentencia}{\nt{If}}
\regla{Sentencia}{\nt{While}}
\regla{Sentencia}{\nt{Bloque}}
\regla{Sentencia}{\nt{For}}\hfill\mbox{\footnotesize\color{grisMedio}[+\,\texttt{for}]}
\end{producciones}

\noindent{\small\color{grisMedio}Sin variables locales clásicas, una sentencia-expresión puede empezar con idClase sin competir con nada: es la regla original.\par}

\begin{producciones}
\regla{Sentencia}{\nt{SentenciaExpresion} \tl{;}}\hfill\mbox{\footnotesize\color{grisMedio}[-\,\texttt{varclasica}]}
\regla{SentenciaExpresion}{\nt{Expresion}}\hfill\mbox{\footnotesize\color{grisMedio}[-\,\texttt{varclasica}]}
\end{producciones}

\noindent{\small\color{grisMedio}Con el logro, \texttt{C x;} y \texttt{C.m();} compiten por idClase y hay que factorizarlo.\par}

\begin{producciones}
\regla{Sentencia}{\nt{ExpresionSinIdClase} \tl{;}}\hfill\mbox{\footnotesize\color{grisMedio}[+\,\texttt{varclasica}]}
\regla{Sentencia}{\nt{TipoPrimitivo} \nt{DimensionesOpcionales} \nt{ListaDeclaradores} \tl{;}}\hfill\mbox{\footnotesize\color{grisMedio}[+\,\texttt{varclasica}]}
\regla{Sentencia}{\tl{idGen} \nt{DimensionesOpcionales} \nt{ListaDeclaradores} \tl{;}}\hfill\mbox{\footnotesize\color{grisMedio}[+\,\texttt{varclasica}]}
\regla{Sentencia}{\tl{idClase} \nt{SentenciaDesdeIdClase}}\hfill\mbox{\footnotesize\color{grisMedio}[+\,\texttt{varclasica}]}
\regla{SentenciaDesdeIdClase}{\nt{TipoGenericoOpcional} \nt{DimensionesOpcionales} \nt{ListaDeclaradores} \tl{;}}\hfill\mbox{\footnotesize\color{grisMedio}[+\,\texttt{varclasica}]}
\regla{SentenciaDesdeIdClase}{\tl{.} \tl{idMetVar} \nt{ArgsActuales} \nt{ColaDeExpresionDesdePrimario} \tl{;}}\hfill\mbox{\footnotesize\color{grisMedio}[+\,\texttt{varclasica}]}
\regla{VarLocal}{\tl{var} \tl{idMetVar} \tl{=} \nt{ExpresionCompuesta}}
\end{producciones}

\noindent{\small\color{grisMedio}Varias variables por sentencia, cada una con inicialización opcional: \texttt{int x,y,z = 10;}\par}

\begin{producciones}
\regla{ListaDeclaradores}{\tl{idMetVar} \nt{InicializacionDeDeclarador} \nt{RestoListaDeclaradores}}\hfill\mbox{\footnotesize\color{grisMedio}[+\,\texttt{varclasica}]}
\regla{RestoListaDeclaradores}{\tl{,} \tl{idMetVar} \nt{InicializacionDeDeclarador} \nt{RestoListaDeclaradores}}\hfill\mbox{\footnotesize\color{grisMedio}[+\,\texttt{varclasica}]}
\regla{RestoListaDeclaradores}{\eps}\hfill\mbox{\footnotesize\color{grisMedio}[+\,\texttt{varclasica}]}
\regla{InicializacionDeDeclarador}{\tl{=} \nt{ExpresionCompuesta} \alt{} \eps}\hfill\mbox{\footnotesize\color{grisMedio}[+\,\texttt{varclasica}]}
\regla{Return}{\tl{return} \nt{ExpresionOpcional}}
\regla{ExpresionOpcional}{\nt{Expresion} \alt{} \eps}
\regla{If}{\tl{if} \tl{(} \nt{Expresion} \tl{)} \nt{Sentencia} \nt{ElseOpcional}}
\regla{ElseOpcional}{\tl{else} \nt{Sentencia} \alt{} \eps}
\regla{While}{\tl{while} \tl{(} \nt{Expresion} \tl{)} \nt{Sentencia}}
\end{producciones}

\noindent{\small\color{grisMedio}[L:for] — Fors! E2 Una sola sentencia/expresión por separador. El \texttt{:} distingue el for de iteradores del estándar, y aparece justo después del declarador, así que el desempate es de un solo token.\par}

\begin{producciones}
\regla{For}{\tl{for} \tl{(} \nt{EncabezadoFor} \tl{)} \nt{Sentencia}}\hfill\mbox{\footnotesize\color{grisMedio}[+\,\texttt{for}]}
\regla{EncabezadoFor}{\tl{;} \nt{ExpresionOpcional} \tl{;} \nt{ExpresionOpcional}}\hfill\mbox{\footnotesize\color{grisMedio}[+\,\texttt{for}]}
\regla{EncabezadoFor}{\tl{var} \tl{idMetVar} \nt{EncabezadoForTrasVar}}\hfill\mbox{\footnotesize\color{grisMedio}[+\,\texttt{for}]}
\regla{EncabezadoFor}{\nt{TipoPrimitivo} \nt{DimensionesOpcionales} \tl{idMetVar} \nt{EncabezadoForTrasDeclarador}}\hfill\mbox{\footnotesize\color{grisMedio}[+\,\texttt{for}]}
\regla{EncabezadoFor}{\tl{idGen} \nt{DimensionesOpcionales} \tl{idMetVar} \nt{EncabezadoForTrasDeclarador}}\hfill\mbox{\footnotesize\color{grisMedio}[+\,\texttt{for}]}
\regla{EncabezadoFor}{\tl{idClase} \nt{EncabezadoForDesdeIdClase}}\hfill\mbox{\footnotesize\color{grisMedio}[+\,\texttt{for}]}
\regla{EncabezadoFor}{\nt{ExpresionSinIdClase} \nt{ClausulasFinalesFor}}\hfill\mbox{\footnotesize\color{grisMedio}[+\,\texttt{for}]}
\regla{EncabezadoForDesdeIdClase}{\nt{TipoGenericoOpcional} \nt{DimensionesOpcionales} \tl{idMetVar} \nt{EncabezadoForTrasDeclarador}}\hfill\mbox{\footnotesize\color{grisMedio}[+\,\texttt{for}]}
\regla{EncabezadoForDesdeIdClase}{\tl{.} \tl{idMetVar} \nt{ArgsActuales} \nt{ColaDeExpresionDesdePrimario} \nt{ClausulasFinalesFor}}\hfill\mbox{\footnotesize\color{grisMedio}[+\,\texttt{for}]}
\end{producciones}

\noindent{\small\color{grisMedio}El \texttt{:} distingue el for de iteradores del estándar, y está a un solo token.\par}

\begin{producciones}
\regla{EncabezadoForTrasVar}{\tl{:} \nt{Expresion}}\hfill\mbox{\footnotesize\color{grisMedio}[+\,\texttt{for}]}
\regla{EncabezadoForTrasVar}{\tl{=} \nt{ExpresionCompuesta} \nt{ClausulasFinalesFor}}\hfill\mbox{\footnotesize\color{grisMedio}[+\,\texttt{for}]}
\regla{EncabezadoForTrasDeclarador}{\tl{:} \nt{Expresion}}\hfill\mbox{\footnotesize\color{grisMedio}[+\,\texttt{for}]}
\regla{EncabezadoForTrasDeclarador}{\nt{InicializacionDeDeclaradorFor} \nt{ClausulasFinalesFor}}\hfill\mbox{\footnotesize\color{grisMedio}[+\,\texttt{for}]}
\regla{InicializacionDeDeclaradorFor}{\tl{=} \nt{ExpresionCompuesta} \alt{} \eps}\hfill\mbox{\footnotesize\color{grisMedio}[+\,\texttt{for}]}
\end{producciones}

\noindent{\small\color{grisMedio}Una sola expresión por separador, como pide el logro.\par}

\begin{producciones}
\regla{ClausulasFinalesFor}{\tl{;} \nt{ExpresionOpcional} \tl{;} \nt{ExpresionOpcional}}\hfill\mbox{\footnotesize\color{grisMedio}[+\,\texttt{for}]}
\end{producciones}

\noindent{\small\color{grisMedio}Capas, de menor a mayor precedencia: Expresion            asignación (NO asociativa, como en la gramática original) ExpresionCondicional ternario [L:ternario] ExpresionCompuesta   operadores binarios (un solo nivel, como el original) ExpresionBasica      operadores unarios prefijos Operando             literal, o referencia con posfijo opcional [L:posfijo]\par}

\begin{producciones}
\regla{Expresion}{\nt{ExpresionCondicional} \nt{RestoExpresion}}
\end{producciones}

\noindent{\small\color{grisMedio}El lado derecho es ExpresionCondicional y NO Expresion: la gramática original pone \nt{ExpresionCompuesta} ahí, o sea que \texttt{a = b = c} NO es válido en MiniJava.\par}

\begin{producciones}
\regla{RestoExpresion}{\nt{OperadorAsignacion} \nt{ExpresionCondicional} \alt{} \eps}
\regla{OperadorAsignacion}{\tl{=}}
\regla{ExpresionCondicional}{\nt{ExpresionCompuesta} \nt{RestoCondicional}}
\end{producciones}

\noindent{\small\color{grisMedio}Con el logro apagado \nt{RestoCondicional} queda en solo ε, o sea que \nt{ExpresionCondicional} se vuelve un alias de \nt{ExpresionCompuesta} y la capa desaparece sin dejar rastro. Asociativo a derecha, como en Java.\par}

\begin{producciones}
\regla{RestoCondicional}{\tl{?} \nt{ExpresionCondicional} \tl{:} \nt{ExpresionCondicional}}\hfill\mbox{\footnotesize\color{grisMedio}[+\,\texttt{ternario}]}
\regla{RestoCondicional}{\eps}
\regla{ExpresionCompuesta}{\nt{ExpresionBasica} \nt{RestoExpresionCompuesta}}
\regla{RestoExpresionCompuesta}{\nt{OperadorBinario} \nt{ExpresionBasica} \nt{RestoExpresionCompuesta} \alt{} \eps}
\regla{OperadorBinario}{\tl{||} \alt{} \tl{\&\&} \alt{} \tl{==} \alt{} \tl{!=} \alt{} \tl{<} \alt{} \tl{>} \alt{} \tl{<=} \alt{} \tl{>=} \alt{} \tl{+} \alt{} \tl{-} \alt{} \tl{*} \alt{} \tl{/} \alt{} \tl{\%}}
\regla{ExpresionBasica}{\nt{OperadorUnario} \nt{Operando}}
\regla{ExpresionBasica}{\nt{Operando}}
\regla{OperadorUnario}{\tl{+} \alt{} \tl{-} \alt{} \tl{!}}
\regla{Operando}{\nt{Primitivo}}
\regla{Operando}{\nt{Referencia} \nt{PosfijoOpcional}}
\end{producciones}

\noindent{\small\color{grisMedio}Va sobre \nt{Referencia} y no sobre \nt{Operando}: \texttt{5++} no tiene sentido.\par}

\begin{producciones}
\regla{PosfijoOpcional}{\tl{++}}\hfill\mbox{\footnotesize\color{grisMedio}[+\,\texttt{posfijo}]}
\regla{PosfijoOpcional}{\tl{-{}-}}\hfill\mbox{\footnotesize\color{grisMedio}[+\,\texttt{posfijo}]}
\regla{PosfijoOpcional}{\eps}
\end{producciones}

\noindent{\small\color{grisMedio}[L:float] — floatLiteral\par}

\begin{producciones}
\regla{Primitivo}{\tl{true} \alt{} \tl{false} \alt{} \tl{intLiteral} \alt{} \tl{charLiteral} \alt{} \tl{null}}
\regla{Primitivo}{\tl{floatLiteral}}\hfill\mbox{\footnotesize\color{grisMedio}[+\,\texttt{float}]}
\end{producciones}

\noindent{\small\color{grisMedio}--- Expresión que no puede empezar con idClase ---------------------------- Se usa donde una expresión compite con una declaración de variable local (sentencia y encabezado del for). Sólo duplica las *cabezas*: todas las colas se comparten con la jerarquía normal, así que no hay lógica repetida.\par}

\begin{producciones}
\regla{ExpresionSinIdClase}{\nt{OperadorUnario} \nt{Operando} \nt{ColaTrasOperando}}
\regla{ExpresionSinIdClase}{\nt{Primitivo} \nt{ColaTrasOperando}}
\regla{ExpresionSinIdClase}{\nt{PrimarioSinIdClase} \nt{ColaDeExpresionDesdePrimario}}
\regla{ColaTrasOperando}{\nt{RestoExpresionCompuesta} \nt{RestoCondicional} \nt{RestoExpresion}}
\regla{ColaDeExpresionDesdePrimario}{\nt{RestoReferencia} \nt{PosfijoOpcional} \nt{ColaTrasOperando}}
\regla{Referencia}{\nt{Primario} \nt{RestoReferencia}}
\regla{RestoReferencia}{\nt{Encadenado} \nt{RestoReferencia}}
\regla{RestoReferencia}{\nt{AccesoArreglo} \nt{RestoReferencia}}
\regla{RestoReferencia}{\eps}
\regla{Encadenado}{\tl{.} \tl{idMetVar} \nt{EncadenadoAtributoOMetodo}}
\end{producciones}

\noindent{\small\color{grisMedio}Dedicado a este contexto: vacío = era <VarEncadenada>, con argumentos = era <MetodoEncadenado>.\par}

\begin{producciones}
\regla{EncadenadoAtributoOMetodo}{\nt{ArgsActuales} \alt{} \eps}
\regla{AccesoArreglo}{\tl{[} \nt{Expresion} \tl{]}}
\regla{Primario}{\nt{LlamadaMetodoEstatico}}
\regla{Primario}{\nt{PrimarioSinIdClase}}
\regla{LlamadaMetodoEstatico}{\tl{idClase} \tl{.} \tl{idMetVar} \nt{ArgsActuales}}
\regla{PrimarioSinIdClase}{\tl{this}}
\regla{PrimarioSinIdClase}{\tl{stringLiteral}}
\regla{PrimarioSinIdClase}{\nt{ParentizadaOLambda}}
\regla{PrimarioSinIdClase}{\nt{AccesoVarLlamadaOLambda}}
\regla{PrimarioSinIdClase}{\nt{Creacion}}
\end{producciones}

\noindent{\small\color{grisMedio}Dedicado a este contexto: vacío = era <AccesoVar>, con argumentos = era <LlamadaMetodo>, con flecha = lambda de un parámetro sin paréntesis.\par}

\begin{producciones}
\regla{AccesoVarLlamadaOLambda}{\tl{idMetVar} \nt{AccesoLlamadaOLambda}}
\regla{AccesoLlamadaOLambda}{\nt{ArgsActuales}}
\regla{AccesoLlamadaOLambda}{\tl{->} \nt{CuerpoLambda}}\hfill\mbox{\footnotesize\color{grisMedio}[+\,\texttt{lambda}]}
\regla{AccesoLlamadaOLambda}{\eps}
\end{producciones}

\noindent{\small\color{grisMedio}[L:lambda] — MiniLambdas E2 \texttt{(} abre tanto una expresión parentizada como la lista de parámetros de una lambda. Parece necesitar lookahead ilimitado (\texttt{(a,b,c) ->} vs \texttt{(a+b)}), pero NO lo necesita: - \texttt{( )}      sólo puede ser una lambda de cero parámetros; - \texttt{( e , …}  sólo puede ser una lambda: MiniJava no tiene operador coma; - \texttt{( e )}    se decide con el token que sigue al \texttt{)}: \texttt{->} es lambda, cualquier otro es una expresión parentizada. Con eso alcanza un token de lookahead en cada punto.\par}

\begin{producciones}
\regla{ParentizadaOLambda}{\tl{(} \nt{ContenidoDeParentesis}}
\regla{ContenidoDeParentesis}{\tl{)} \tl{->} \nt{CuerpoLambda}}\hfill\mbox{\footnotesize\color{grisMedio}[+\,\texttt{lambda}]}
\regla{ContenidoDeParentesis}{\nt{Expresion} \nt{RestoDeParentesis}}
\regla{RestoDeParentesis}{\tl{)} \nt{FinDeParentizadaOLambda}}
\regla{RestoDeParentesis}{\tl{,} \nt{ListaParametrosLambda} \tl{)} \tl{->} \nt{CuerpoLambda}}\hfill\mbox{\footnotesize\color{grisMedio}[+\,\texttt{lambda}]}
\regla{FinDeParentizadaOLambda}{\tl{->} \nt{CuerpoLambda}}\hfill\mbox{\footnotesize\color{grisMedio}[+\,\texttt{lambda}]}
\regla{FinDeParentizadaOLambda}{\eps}
\regla{ListaParametrosLambda}{\tl{idMetVar} \nt{RestoListaParametrosLambda}}\hfill\mbox{\footnotesize\color{grisMedio}[+\,\texttt{lambda}]}
\regla{RestoListaParametrosLambda}{\tl{,} \tl{idMetVar} \nt{RestoListaParametrosLambda} \alt{} \eps}\hfill\mbox{\footnotesize\color{grisMedio}[+\,\texttt{lambda}]}
\end{producciones}

\noindent{\small\color{grisMedio}Una única expresión, sin llaves.\par}

\begin{producciones}
\regla{CuerpoLambda}{\nt{Expresion}}\hfill\mbox{\footnotesize\color{grisMedio}[+\,\texttt{lambda}]}
\regla{Creacion}{\tl{new} \nt{ArregloOInstancia}}
\regla{ArregloOInstancia}{\nt{TipoPrimitivo} \nt{CreacionDeArreglo}}
\regla{ArregloOInstancia}{\tl{idGen} \nt{CreacionDeArreglo}}
\regla{ArregloOInstancia}{\tl{idClase} \nt{CreacionDesdeIdClase}}
\regla{CreacionDesdeIdClase}{\tl{<} \nt{GenericoDeCreacion}}
\regla{CreacionDesdeIdClase}{\nt{CreacionDeArreglo}}
\regla{CreacionDesdeIdClase}{\nt{ArgsActuales}}
\end{producciones}

\noindent{\small\color{grisMedio}[L:generico] — el diamante \texttt{<>} sólo se admite acá, en la rama que va a ArgsActuales: \texttt{new C<>()} es válido y \texttt{new C<>[3]} no. La restricción del logro queda impuesta por la gramática, no delegada al semántico.\par}

\begin{producciones}
\regla{GenericoDeCreacion}{\tl{>} \nt{ArgsActuales}}\hfill\mbox{\footnotesize\color{grisMedio}[+\,\texttt{generico}]}
\regla{GenericoDeCreacion}{\nt{ArgumentoDeTipo} \tl{>} \nt{ArregloOInstanciaDeClase}}
\regla{ArregloOInstanciaDeClase}{\nt{CreacionDeArreglo}}
\regla{ArregloOInstanciaDeClase}{\nt{ArgsActuales}}
\end{producciones}

\noindent{\small\color{grisMedio}[L:arrinit] — \texttt{new int[]\{1,2,3\}}. Sin el logro: CreacionDeArreglo -> [ Expresion ] DimensionesConTamanioOpcional\par}

\begin{producciones}
\regla{CreacionDeArreglo}{\tl{[} \nt{RestoCreacionDeArreglo}}
\regla{RestoCreacionDeArreglo}{\nt{Expresion} \tl{]} \nt{DimensionesConTamanioOpcional}}
\regla{RestoCreacionDeArreglo}{\tl{]} \tl{\{} \nt{ListaExpsOpcional} \tl{\}}}\hfill\mbox{\footnotesize\color{grisMedio}[+\,\texttt{arrinit}]}
\regla{DimensionesConTamanio}{\tl{[} \nt{Expresion} \tl{]} \nt{DimensionesConTamanioOpcional}}
\regla{DimensionesConTamanioOpcional}{\nt{DimensionesConTamanio} \alt{} \eps}
\regla{ArgsActuales}{\tl{(} \nt{ListaExpsOpcional} \tl{)}}
\regla{ListaExpsOpcional}{\nt{ListaExps} \alt{} \eps}
\regla{ListaExps}{\nt{Expresion} \nt{RestoListaExps}}
\regla{RestoListaExps}{\tl{,} \nt{Expresion} \nt{RestoListaExps} \alt{} \eps}
\end{producciones}

